\let\negmedspace\undefined
\let\negthickspace\undefined
\documentclass[journal,onecolumn]{IEEEtran}

\usepackage[a5paper, margin=10mm]{geometry} % Safe with onecolumn
\usepackage{lmodern} % Better fonts (optional but recommended)
\usepackage{tfrupee} % Rupee symbol package

% Header fix
\setlength{\headheight}{12pt} % Safe value
\setlength{\headsep}{0pt}     % Reduce gap

\usepackage{gvv-book}
\usepackage{gvv}
\usepackage{cite}
\usepackage{amsmath,amssymb,amsfonts,amsthm}
\usepackage{algorithmic}
\usepackage{graphicx}
\usepackage{textcomp}
\usepackage{xcolor}
\usepackage{txfonts}
\usepackage{listings}
\usepackage{enumitem}
\usepackage{mathtools}
\usepackage{gensymb}
\usepackage{comment}
\usepackage[breaklinks=true]{hyperref}
\usepackage{tkz-euclide} 
\usepackage{listings}
% \usepackage{gvv}                                        
\def\inputGnumericTable{}                                 
\usepackage[latin1]{inputenc}                                
\usepackage{color}                                            
\usepackage{array}                                            
\usepackage{longtable}                                       
\usepackage{calc}                                             
\usepackage{multirow}                                         
\usepackage{hhline}                                           
\usepackage{ifthen}                                           
\usepackage{lscape}
\usepackage{multicol}

\usepackage{amsmath,amssymb}
\usepackage{booktabs}
\usepackage{tikz}
\usetikzlibrary{arrows.meta,angles,quotes}

\begin{document}

\bibliographystyle{IEEEtran}
\vspace{3cm}

\title{5.13.40}
\author{EE25BTECH11014 - Bhoomika Lokesh}
% \maketitle
% \newpage
% \bigskip
{\let\newpage\relax\maketitle}

\renewcommand{\thefigure}{\theenumi}
\renewcommand{\thetable}{\theenumi}
\setlength{\intextsep}{10pt} % Space between text and floats

% Remove this line so equations use default numbering (1, 2, 3,...)
% \numberwithin{equation}{enumi}
\numberwithin{figure}{enumi}
\renewcommand{\thetable}{\theenumi}


\textbf{Question:} Let $a,\lambda,\mu \in \mathbb{R}$. Consider the system of linear equations\\
\begin{center}
$ax+2y=\lambda$\\
$3x-2y=\mu$
\end{center}
Which of the following statement$\brak{s}$ is $\brak{are}$ correct?
\begin{multicols}{1}
\begin{enumerate}
\item If $a=-3,$ then the system has infinitely many solutions for all value of $\lambda$ and $\mu$.
 \item If $a \neq -3$, then the system has unique solution for all values of $\lambda$ and $\mu.$
 \item  If $\lambda + \mu = 0$, then the system has infinitely many solutions for $a = -3$.
\item  If $\lambda + \mu \neq 0$, then the system has no solution for $a = -3$.
\end{enumerate}
\end{multicols}

\textbf{Solution:}\\
By the system of equations,we get the following augmented matrix:\\
By \ row \ reduction \ method,
\begin{align} 
\myvec{\begin{array}{cc|c}
a & 2 & \lambda \\
3 & -2 & \mu
\end{array}} \xrightarrow{R_2\longrightarrow R_2-3\frac{R_1}{a}}\myvec{\begin{array}{cc|c}
a & 2 & \lambda \\
0 & \frac{-2a-6}{a} & \mu-\frac{3\lambda}{a}
\end{array}} \\
Now \ if \ \frac{-2a-6}{a}=0 \implies a=-3,\\
when \ a=-3,\mu-3\frac{\lambda}{a}=\mu+\lambda
\end{align}

\textbf{Case-1:}If $\lambda+\mu=0$ and the system has \textbf{infinitely many solutions}.\\
\textbf{Case-2:}If $\lambda+\mu\neq0$, then the system has \textbf{no solution}\\
\\and if $a\neq-3$, then the first row is non-zero and matrix is invertible. Hence system has a \textbf{unique solution} for all values of $\lambda$ and $\mu$. \\

And solution for $a\neq-3$ is:
\begin{align}
   x\brak{a+3}=\lambda+\mu \implies x=\frac{\lambda+\mu}{a+3}\\
   3x-2y=\mu \implies y=\frac{3\lambda-a\mu}{2a+6}
 \end{align}

Therefore (2),(3) and (4) are the correct statements.

\end{document}
